% !TeX root = ../papers/nc-p3.tex
\section{階数15の全通常窓の二次双対}
\label{sec:rank-fifteen-quadratic-dual}
\begin{theorem}\label{thm:rank-fifteen-quadratic-dual}
定理\ref{thm:rank-fifteen-covers}の両表示の全通常な階数15窓で、各始点の二次双対は
次数次元 $(1,4,6,4,1)$ の完全対合を持ち、次数5以上は零である。
特に前節に残る孤立点でも有限Frobenius条件は成立し、標準Koszul左右AS正則性は
元の代数のKoszul性と同値である。
\end{theorem}
\begin{proof}
位置 $p$ の二次関係を $R_p$、長さ $d$ の全ての隣接位置に $R$ を置いた空間の
共通部分を $K_{p,d}$ とする。これが二次双対の長さ $d$ の双対である。
文字の和を $\oplus$ と書く。全文字 $h$ の三次テンソルを、最初と最後の文字を
行 $i$、列 $k$ とする四行四列の表に書く。中央の文字は $i\oplus h\oplus k$ で一意である。
$K_{p,3}$ では列 $k$ が $R_p[h\oplus k]$ に、行 $i$ が
$R_{p+1}[h\oplus i]$ に入る。行の座標順は中央の文字と最後の文字の変換で並べ替える。

$J=h\oplus\{2,3\}$ に属する二行と二列のそれぞれは、全成分が非零の既知ベクトルの
スカラー倍である。$J\times J$ での四つの一致式により、この四本の倍率は
一つの倍率から決まる。残る行の関係空間は二次元で、対応する文字 $g$ は0または1である。
$J$ は $\{0,1\}$ または $\{2,3\}$ であり、$g\oplus J=J$ である。
厚い二次関係を $(L\mid I_2)$ と書くと $L$ も可逆なので、
$J$ 内の二成分を指定すれば残る二成分まで決まる。従って三次交叉は文字ごとに
高々一次元である。\path{computations/verify_rank_fifteen_quadratic_dual.py} は
四位置の $\det L$ の分子・分母が通常な単元の積であり、細い関係の全成分も単元である
ことを斉次座標から検算する。四階テンソル $W_p$ の一文字縮約が
各文字に非零の元を与えるので、これらがちょうど $K_{p,3}$ を張る。

$K_{p,4}$ の元は両端からの三次縮約により
\[
 (A\otimes1\otimes1\otimes1)W_p
   =(1\otimes1\otimes1\otimes B)W_p
\]
と書ける。反対切断を第一・第三文字を行、第二・第四文字を列とする行列 $M$ にすれば、
$(A\otimes1)M=M(1\otimes B^{\mathsf T})$ となる。
階数15では $M$ の核が一本であり、その生成元は二文字間の可逆な四行四列行列である。
実際、定理\ref{thm:fixed-opposite-rank-drop}のブロック分解で核を持つ文字は0または1で、
対反対称部分は可逆である。単元で行または列を規格化した対称ブロックは、
対称座標 $(t,t,s,s)$ 上で対角外成分がともに2の二行二列行列になる。
従って核の非零ベクトルは $t,s$ ともに非零である。これを二文字の行列に戻すと
$j$ 行 $j\oplus g$ 列に一つずつ非零成分があり、確かに可逆である。
規格化したブロックは対称なので左核にも同じ結論が成り立ち、
位置を巡回してもこの性質は保たれる。

$1\otimes B^{\mathsf T}$ はこの核直線を保つ。可逆な二文字行列に片側から作用して
スカラー倍となるため $B$ 自身がスカラーである。よって $K_{p,4}=kW_p$ である。
五次の元は $W_p\otimes v=w\otimes W_{p+1}$ と書けるが、非零なら第一文字での
$W_p$ の階数が1になり、階数4に反する。これで五次以後も零となる。
$W_p$ の一文字切断と隣接二文字切断の階数4,6により、最高次数への積は次数
$(1,3),(2,2),(3,1)$ で完全対合を与える。Koszulならこの二次双対から左右の
長さ4の自己双対な解消が得られ、Hilbert級数は $(1-t)^{-4}$ となるので主張が従う。
\end{proof}

\begin{corollary}\label{cor:rank-fifteen-sixth-test}
これらの窓のHilbert次元が最初に標準値を超え得るのは長さ6である。
各始点で文字0から3の順に、末尾の二次関係の挿入行列の大きさは
$56\times54,56\times54,56\times51,56\times51$ である。
長さ6の次元が84を超える条件は、このいずれかがそれぞれの目標階数
$36,36,34,34$ を下回ることと同値である。
\end{corollary}
\begin{proof}
長さ4,5は定理\ref{thm:fixed-opposite-rank-drop}により全通常点で一定であり、
各文字では一般点の下界に等しく $(9,10,8,8),(14,14,14,14)$ である。
長さ3は上の交叉から各文字で $16-6-6+1=5$ となる。
長さ6の商は $A_4\otimes R\longrightarrow A_5\otimes V$ の余核である
（各因子は対応する位置で取る）。出力は文字ごとに56次元、入力は関係の文字次元
$(2,2,1,1)$ と $A_4$ の畳み込みで上記の大きさを持つ。
正則な一般窓では直前の二項の次元がそれぞれ20と
$\dim A_2[h]=(2,2,3,3)_h$ である。従って完全性から表示した四階数を得る。
一定階数の経路基底の上で行列の階数は特殊化により増えず、目標時の余核次元の合計は
$20+20+22+22=84$ である。上記コードは各位置でこの挿入行列を直接生成する。
$\xi=9/16$、101上の $K=2$ では $H^2=3$ が非平方である。
二次剰余体で全八シートと両表示の行列を実際に構成し、全ての目標階数を確認した。
\end{proof}
